\section{Campagne completate e tracciabilità}\label{app:campaigns}
La Tabella~\ref{tab:campaigns} distingue gli esperimenti storici dagli approfondimenti disponibili con la revisione. I conteggi si riferiscono a run, record o lotti secondo l'unità indicata. Le configurazioni estese presenti nel software costituiscono proposte di ulteriori esperimenti e sono escluse dai risultati di questo articolo.
\begin{table}[H]\centering\small
\caption{Inventario delle evidenze usate nel manoscritto.}\label{tab:campaigns}
\begin{tabularx}{\textwidth}{@{}p{3.3cm}Xr@{}}\toprule
Esperimento & Disegno completato & Quantità\\\midrule
Screening del passo & 6 funzioni, 2 dimensioni, 5 passi, 5 semi, $K=30$ & 300 run\\
Validazione del passo & 12 celle, 4 passi, 30 nuovi semi, $K=30$ & 1\,440 run\\
Topologie interne & 4 celle, 8 grafi, 24 semi, 3 ripetizioni, 4 worker count & 9\,216 run\\
BBOB iniziale & 24 funzioni, 2 dimensioni, 5 istanze, 12 metodi & 2\,880 run\\
Calibrazione BBOB & 24 funzioni, 2 dimensioni, 3 istanze, 4 passi & 576 run\\
Validazione BBOB & 24 funzioni, 2 dimensioni, 5 istanze, 13 etichette & 3\,120 record\\
CEC a budget ridotto & 15 funzioni, 5 semi, 12 metodi & 900 run\\
Replica BBOB $D=5$ & 24 funzioni, 5 istanze, 11 metodi, nuovo seme & 1\,320 run\\
Ablazione controllata & 4 funzioni, 2 dimensioni, 3 istanze, 65 metodi & 1\,560 run\\
Costo e parallelismo & 3 sessioni, 4 costi, 2 dinamiche, 2 ripetizioni, 12 modalità; 8 semi per lotto & 576 lotti\\
Diagnostica affine & 3 dimensioni, 3 semi, 16 configurazioni & 144 run\\
\bottomrule\end{tabularx}\end{table}

\subsection{Parametri condivisi ed eccezioni}
Il nucleo usa $c_c=c_s=2.05$, $\mu_{\max}=0.9$, $\mu_{\min}=0.4$, velocità iniziale uniforme in $[-0.1w_j,0.1w_j]$ e restituzione $e=0.5$. La popolazione è $N=64$. Il passo delle topologie e degli approfondimenti temporali e geometrici è $1$. Il fattoriale varia il passo analitico; la comparazione esterna trasferisce quello selezionato nel proprio disegno. Le schedule d'inerzia dipendono dal budget programmato: un checkpoint intermedio appartiene alla traiettoria del budget massimo.

\subsection{Versioni e identificazione dei dati}
La base del manoscritto e il codice consultato corrispondono al commit:\par\noindent
\texttt{c56be76eda7dd7299bd697928576b0bb286b6297}.
I metadati delle campagne recenti registrano come HEAD:\par\noindent
\texttt{bbdedd02346cb62209aa596ffb8a076d01ee2ae3}; gli hash dei singoli sorgenti e del binario identificano anche eventuali modifiche locali rispetto a quell'HEAD. La build dichiarata è AppleClang 21.0.0.21000099, C++20, Release, arm64, opzioni \texttt{-O3 -DNDEBUG}, ABI 7. I risultati provengono dalla macchina dell'autore; la compilazione del manoscritto è un'operazione distinta.

I riepiloghi recenti allegati ai sorgenti contengono i disegni, la provenienza, i ranghi, tutti i contrasti, le ECDF, gli ERT e i timer di fase. La famiglia BBOB controllata utilizza esclusivamente la versione corretta con minimo letto dall'Observer dopo la chiusura del logger. Gli output delle esecuzioni precedenti interrotte sono esclusi. L'analisi statistica successiva ricava le tre matrici dalla validazione storica, con undici metodi: l'esclusione dei controlli CPSO aggiuntivi modifica i ranghi e la famiglia di Holm rispetto alle tabelle originarie.

\subsection{Graduatorie visive della campagna topologica interna}
La Figura~\ref{fig:restored-topology-ranks} riproduce la graduatoria completa di qualità e tempi della campagna interna a $h=1$. Le colonne misurano criteri diversi: rango della fitness finale e rango del tempo con quattro worker su Rastrigin e otto sulla diffusione inversa. Il confronto va distinto dal fattoriale BBOB della successiva Appendice~\ref{app:controlled-results}.
\begin{figure}[H]\centering
\includegraphics[width=0.90\textwidth,height=0.76\textheight,keepaspectratio]{fig_topology_ranks.pdf}
\caption{Ranghi appaiati delle otto topologie nelle quattro celle obiettivo--dimensione: qualità finale (sinistra) e tempo (destra: quattro worker su Rastrigin, otto su diffusione inversa). Le barre riportano gli intervalli della sintesi originaria per seme. ``Wheel (star)'' identifica la stella.}\label{fig:restored-topology-ranks}
\end{figure}


\subsection{Ambiente temporale osservato}
I metadati delle tre sessioni registrano otto CPU logiche e quattro core in ciascuno dei due livelli di prestazione M2. Lo stato di alimentazione osservato è collegato alla rete elettrica, con batteria carica. I campi relativi a modalità a basso consumo, temperatura e stato termico sono non disponibili. Le sessioni 1--2 addestrano il modello e la 3 verifica la scelta, sulla stessa macchina. La calibrazione del lavoro viene ripetuta per sessione; i confronti di equivalenza raggruppano per il lavoro intero effettivamente eseguito.

\subsection{Regole di sintesi}
Nel confronto storico delle topologie si aggregano prima le ripetizioni di uno stesso seme. Nella replica dei grafi si aggregano le realizzazioni entro sciame. Nell'analisi della famiglia PSO il bootstrap ricampiona le funzioni conservando le cinque istanze; il test dei segni usa le mediane delle differenze appaiate per funzione. Nel fattoriale le tre istanze sono aggregate entro funzione prima dell'inferenza. Nel costo controllato la tabella temporale riporta la mediana dei rapporti delle due ripetizioni nella terza sessione: i suoi valori sono descrittivi, mentre i lotti contengono gli stessi otto semi.


\FloatBarrier
